\documentclass[11pt]{article}
\usepackage{homework}

\course{Math 113: Abstract Algebra}
\term{Fall 2026}
\title{Homework 11}
\author{Jonathan Nguyen}

\begin{document}

\begin{problem}[26.4]
    Give addition and multiplication tables for $2\mathbb{Z}/8\mathbb{Z}$. Are $2\mathbb{Z}/8\mathbb{Z}$ and $\mathbb{Z}_4$ isomorphic?
\end{problem}
\begin{solution}
    \[
\begin{array}{c|cccc}
+ & \overline{0} & \overline{2} & \overline{4} & \overline{6} \\
\hline
\overline{0} & \overline{0} & \overline{2} & \overline{4} & \overline{6} \\
\overline{2} & \overline{2} & \overline{4} & \overline{6} & \overline{0} \\
\overline{4} & \overline{4} & \overline{6} & \overline{0} & \overline{2} \\
\overline{6} & \overline{6} & \overline{0} & \overline{2} & \overline{4}
\end{array}
\qquad
\begin{array}{c|cccc}
\cdot & \overline{0} & \overline{2} & \overline{4} & \overline{6} \\
\hline
\overline{0} & \overline{0} & \overline{0} & \overline{0} & \overline{0} \\
\overline{2} & \overline{0} & \overline{4} & \overline{0} & \overline{4} \\
\overline{4} & \overline{0} & \overline{0} & \overline{0} & \overline{0} \\
\overline{6} & \overline{0} & \overline{4} & \overline{0} & \overline{4}
\end{array}
\]
    From the multiplication table, we know that the products of elements in
    $2\mathbb{Z}/8\mathbb{Z}$ can only be $\overline{0}$ or $\overline{4}$.
    In $\mathbb{Z}_4$, we can get all four elements as products by
    multiplying each element by $1$. Therefore there cannot exist a
    bijection that preserves multiplication, and in turn
$2\mathbb{Z}/8\mathbb{Z} \not\cong \mathbb{Z}_4$.
\end{solution}
\begin{problem}[26.18]
    Show that each homomorphism from a field to a ring is either one to one or maps everything onto 0.
\end{problem}
\begin{solution}
    Let $\phi$ be a ring homomorphism from $F$ to a ring, and let
    $I = \ker(\phi)$, which is an ideal. If $I$ is the trivial ideal
    $\{0\}$, then we know $\phi$ is injective.

    Let $0 \neq a \in I$. By defintion, we know that $a^{-1} \in F$.
    Since, we must satisfy that $a^{-1}I \subseteq I$, we can conclude
    $1 \in I$. Since $1 \in I$, we know that all elements in $F$ are
    also in $I$. Therefore, $I = F$ and $\phi$ maps every element of
    $F$ to $0$.
\end{solution}

\begin{problem}[22.20]
    Let $R$ be a commutative ring with unity of prime characteristic $p$.
    Show that the map $\phi_p: R \rightarrow R$ given by
    $\phi_p(a) = a^p$ is a homomorphism (the \textbf{Frobenius homomorphism}).
\end{problem}
\begin{solution}
    \[
    \phi_p(ab) = (ab)^p = a^p b^p = \phi_p(a)\phi_p(b)
    \]
    We now want to show that $\phi_p(a + b) = \phi_p(a) + \phi_p(b)$.
    \begin{align*}
        \phi_p(a + b) &= (a+b)^p \\
        &= C(p,0)a^p + C(p,1)a^{p-1}b + \cdots + C(p, p - 1)ab^{p-1} + C(p,p)b^p
    \end{align*}
    Since $C(p,k) = \frac{p!}{k!(p-k)!}$ and $p$ divides $p!$, but not $k!$ or $(p-k)!$ when $0 < k < p$. Thus $C(p,k) \equiv 0 \pmod{p}$ when $k \neq 0, p$, so it does not contribute to the binomial expansion. Since $p$ is the characteristic, any element multiplied by $p$ is 0. Therfore, 
    \begin{align*}
        (a+b)^p &= a^p + b^p\\
        &= \phi_p(a) + \phi_p(b)
    \end{align*}
    Therefore $\phi_p$ is a homomorphism.
\end{solution}
\begin{problem}[26.30]
    An element $a$ of a ring $R$ is nilpotent if $a^n = 0$ for some $n \in \mathbb{Z}^+$. Show that the collection of nilpotent elements in a commutative ring $R$ is an ideal, the \textbf{nilradical of} $R$.
    
\end{problem}
\begin{solution}
    We must show that the collection of nilpotent elements in the commutative ring $R$ forms an additive subgroup such that $aI \subseteq I$ for all $a \in R$.

    Let $a \in I$ be a an element from the nilpotent set $I$ where $a^n = 0$, and $b$ an element from $R$.
    \begin{align*}
        (ab)^n = a^nb^n = 0(b^n) = 0
    \end{align*}
    Therefore multiplying any element in the ring $R$ with a nilpotent element is also nilpotent, meaning that $aI \subseteq I$ for all $a \in R$.

    Now we must show that if $a,b \in I$, then $a + b \in I$. Let $l = n + k$ where $a^n = 0$ and $b^k = 0$. We can use binomial expanision to express $(a+b)^l = \sum_{r = 0}^{l}C(l, r)a^{l-r}b^r$. For the terms where $l-r \geq n$, we can rewrite the term as 
    \begin{align*}
        C(l, r)a^{l-r}b^r = C(l,r)a^na^{l-r-n}b^r = C(l,r)(0)a^{l-r-n}b^r = 0
    \end{align*}
    When $l-r < n$, we know that $r > k$ so we can rewrite those terms as 
    \begin{align*}
        C(l, r)a^{l-r}b^r = C(l,r)a^{l-r}b^kb^{r-k} = C(l,r)a^{l-r}(0)b^{r-k} = 0
    \end{align*}
    Since all terms evaluate to 0, $(a+b)^l=0$ and therefore, $a+b \in I$.
    
    Let $a$ be nilpotent where $a^n = 0$. $(-a)^n = (-1)^n(a)^n = (-1)^n(0) = 0$. Therefore $I$ contains additive inverses. $0^1 = 0$, meaning that the additive identity is also in $I$.

    Thus, the nilpotent elements from a commutative ring $R$ form an ideal.
\end{solution}
\end{document}